% from main.tex
%%%%%%%%%%  scan idx 97  |  folio 84  %%%%%%%%%%
covariance de \(K^\bullet\), contravariance de \(K_\bullet\), pour des morphismes satisfaisant des hypothèses de régularité convenables. Enfin, les critères de globalisation de l'Exp.\ II permettent de faire le lien avec les «~groupes de Grothendieck~» définis naïvement à partir des faisceaux cohérents ou des faisceaux localement libres de type fini.

\section*{1. Définitions préliminaires}

\subsection*{1.1. Catégories fibrées à fibres additives (resp. abéliennes, resp. triangulées)}

\textbf{1.1.1.} Soit \(S\) une catégorie, et soit \(\mathcal C\) une \(S\)-catégorie fibrée (SGA 1 VI 6.1) à fibres des catégories additives (resp. triangulées). Nous dirons que \(\mathcal C\) est une \emph{\(S\)-catégorie additive} (\emph{resp.\ triangulée}) (ou que \(\mathcal C\) est additive (resp.\ triangulée) sur \(S\)) si, pour toute flèche \(f:X\to Y\) de \(S\), le foncteur image inverse (déterminé à isomorphisme unique près) \(f^*:\mathcal C_Y\longrightarrow \mathcal C_X\) est additif (resp.\ exact). Nous dirons que \(\mathcal C\) est une \(S\)-catégorie \emph{abélienne} si \(\mathcal C\) est une \(S\)-catégorie additive à fibres abéliennes, et que \(\mathcal C\) est une \(S\)-catégorie abélienne \emph{plate} si en outre les foncteurs image inverse sont exacts.

\textbf{1.1.2.} Soit \(\mathcal C\) une \(S\)-catégorie triangulée. On définit, par exemple à l'aide d'un clivage de \(\mathcal C\) (SGA 1 VI 7.1), une \(S\)-catégorie fibrée \(\Trc(\mathcal C)\) dont la fibre en chaque \(X\in \ob(S)\) est la catégorie \(\Trc(\mathcal C_X)\) des triangles distingués de \(\mathcal C_X\).

\textbf{1.1.3.} Soient \(\mathcal C\) et \(\mathcal C'\) des \(S\)-catégories additives (resp. triangulées) et \(F:\mathcal C\to\mathcal C'\) un foncteur. On dit que \(F\) est un \emph{\(S\)-foncteur additif} (\emph{resp.\ exact}) si \(F\) est un \(S\)-foncteur cartésien additif (resp.\ exact) sur chaque fibre.



% Additional macro definitions
\providecommand{\cC}{\mathcal C}
\providecommand{\cD}{\mathcal D}
\providecommand{\cK}{\mathcal K}
\providecommand{\cO}{\mathcal O}
\providecommand{\cS}{\mathcal S}
\providecommand{\Cpx}{\mathrm C}
\providecommand{\Kpx}{\mathrm K}
\providecommand{\Dpx}{\mathrm D}
\providecommand{\D}{\mathrm D}
\providecommand{\Ho}{\mathrm H}
\providecommand{\RHom}{\mathrm R\!\operatorname{Hom}}
\providecommand{\Ob}{\operatorname{ob}}
\providecommand{\Coker}{\operatorname{Coker}}
\providecommand{\Ker}{\operatorname{Ker}}
\providecommand{\Im}{\operatorname{Im}}
\providecommand{\Hom}{\operatorname{Hom}}
\providecommand{\Cone}{\operatorname{Cone}}
\providecommand{\id}{\operatorname{Id}}
\providecommand{\qis}{\xrightarrow{\sim}}


